\chapter{Die Haarlocke}
\quellseite{9}

„Seht Ihr diese Haarlocke?“ fragte mich ein Greis, indem er mir eine solche zeigte.

Da die Haare blond waren, äußerte ich die Vermuthung, sie rühre von einem Kinde her, das ihm früh weggestorben sei.

„Nein“, antwortete er, „es sind Haare von mir. Es sind jetzt beinahe siebzig Jahre, daß sie von meinem Kopfe abgeschnitten worden sind!“

„So?“ fragte ich, indem mein Interesse erwachte. „Warum habt Ihr sie aufbewahrt?“

„Warum ich das gethan habe?“

„Das will ich Euch erzählen“, sagte der Greis. „Es geschieht zu Gottes Ehre! Ich war höchstens vier Jahre alt — meine langen, goldgelben Haare umspielten meine Wangen und Schultern in glänzenden Locken —, als eines Tages mein Vater in den benachbarten Wald ging, um einen großen Baum zu fällen, wobei ich ihn begleitete. Während der Vater, welcher sehr stark war, mit seiner Art so kräftig in den Stamm schlug, daß die Späne bis zu mir flogen, belustigte ich mich damit, diese zu sammeln. Als ich mich aber ein wenig näherte, um einen Span zu ergreifen, stolperte ich und fiel gerade an den Ort, wo im gleichen Augenblick die Art mit aller Kraft niederfuhr. Der Vater, von Schrecken ergriffen und nicht zweifelnd, daß er mich getödtet habe, konnte nichts thun, als einen Schrei der Verzweiflung ausstoßen. Aber Gott selbst, dieser treue Vater, der die Haare seiner Kinder zählt, hat mich behütet. Die Art hatte im Niederfahren nur diese Locke abgeschnitten, sonst aber meinen Kopf nicht berührt. Als mein Vater die wunderbare Rettung sah, brach er in Thränen aus, nahm mich in seine Arme, fiel auf die Knie nieder, lobte Gott mit lauter Stimme und bat ihn, daß er mich doch zu seinem Ruhme aufwachsen lassen wolle.

„Seit jenem Vorfall hing diese Locke stets im großen Zimmer unter dem Spiegel, und von Zeit zu Zeit zeigte sie mir der Vater, indem er sagte: ‚Mein Sohn, denke jeden Tag daran, daß, wenn die Art, welche auf dich fallen wollte, abgeleitet wurde, du es Gott allein zu verdanken hast! Aber denke auch daran, daß er in seiner Barmherzigkeit deine Seele vor einem viel schrecklicheren Schlage verschont hat, indem er dich durch die Gnade Jesu Christi von der Verdammniß befreit hat!‘“
